\documentclass[12pt,a4paper]{article}

% ─── Packages ──────────────────────────────────────────────────────────────────
\usepackage[utf8]{inputenc}
\usepackage[T1]{fontenc}
\usepackage[french]{babel}
\usepackage{geometry}
\usepackage{graphicx}
\usepackage{booktabs}
\usepackage{array}
\usepackage{xcolor}
\usepackage{hyperref}
\usepackage{amsmath}
\usepackage{amssymb}
\usepackage{enumitem}
\usepackage{fancyhdr}
\usepackage{titlesec}
\usepackage{caption}
\usepackage{subcaption}
\usepackage{float}
\usepackage{multirow}
\usepackage{colortbl}
\usepackage{mdframed}
\usepackage{listings}
\usepackage{tcolorbox}
\usepackage{pgfplots}
\pgfplotsset{compat=1.18}
\usepackage{tikz}

\usepackage{setspace}
\usepackage{abstract}

% ─── Geometry ──────────────────────────────────────────────────────────────────
\geometry{
  top=2.5cm, bottom=2.5cm,
  left=2.8cm, right=2.8cm,
  headheight=15pt
}

% ─── Colors ────────────────────────────────────────────────────────────────────
\definecolor{primary}{RGB}{33,150,243}
\definecolor{secondary}{RGB}{76,175,80}
\definecolor{accent}{RGB}{255,152,0}
\definecolor{danger}{RGB}{244,67,54}
\definecolor{lightgray}{RGB}{245,245,245}
\definecolor{darkgray}{RGB}{66,66,66}
\definecolor{best}{RGB}{200,230,201}

% ─── Header/Footer ─────────────────────────────────────────────────────────────
\pagestyle{fancy}
\fancyhf{}
\fancyhead[L]{\small\textcolor{primary}{\textbf{Prédiction de Retards de Vols}}}
\fancyhead[R]{\small\textcolor{darkgray}{Machine Learning | 2024}}
\fancyfoot[C]{\thepage}
\renewcommand{\headrulewidth}{0.4pt}

% ─── Section styles ────────────────────────────────────────────────────────────
\titleformat{\section}{\Large\bfseries\color{primary}}{}{0em}{}[\titlerule]
\titleformat{\subsection}{\large\bfseries\color{darkgray}}{}{0em}{}
\titleformat{\subsubsection}{\normalsize\bfseries\color{darkgray}}{}{0em}{}

% ─── Hyperref ──────────────────────────────────────────────────────────────────
\hypersetup{
  colorlinks=true, linkcolor=primary, urlcolor=primary, citecolor=secondary,
  pdftitle={Prédiction de Retards de Vols},
  pdfauthor={Projet Machine Learning}
}

% ─── tcolorbox styles ──────────────────────────────────────────────────────────
\tcbuselibrary{skins,breakable}
\newtcolorbox{infobox}[1][]{
  enhanced, colback=primary!8, colframe=primary, fonttitle=\bfseries,
  title=#1, breakable
}
\newtcolorbox{resultbox}[1][]{
  enhanced, colback=secondary!8, colframe=secondary, fonttitle=\bfseries,
  title=#1, breakable
}
\newtcolorbox{warningbox}[1][]{
  enhanced, colback=accent!10, colframe=accent, fonttitle=\bfseries,
  title=#1, breakable
}

% ─── Listings ──────────────────────────────────────────────────────────────────
\lstset{
  language=Python, basicstyle=\small\ttfamily,
  backgroundcolor=\color{lightgray}, frame=single,
  keywordstyle=\color{primary}\bfseries,
  commentstyle=\color{secondary}\itshape,
  stringstyle=\color{danger},
  breaklines=true, showstringspaces=false,
  numbers=left, numberstyle=\tiny\color{darkgray}
}

% ──────────────────────────────────────────────────────────────────────────────
\begin{document}

% ─── Title Page ───────────────────────────────────────────────────────────────
\begin{titlepage}
  \centering
  \vspace*{1.5cm}

  {\Huge\bfseries\color{primary} Prédiction des Retards de Vols}\\[0.5cm]
  {\Large\color{darkgray} Classification Binaire par Machine Learning}\\[2cm]

  \begin{tcolorbox}[enhanced, colback=primary!10, colframe=primary, width=0.8\textwidth]
    \centering
    \textbf{Résumé du Projet}\\[0.3cm]
    \small
    Ce rapport présente une étude complète de prédiction des retards de vols aériens
    à partir d'un dataset de plus de 3 millions d'enregistrements. Quatre modèles de 
    Machine Learning sont comparés : Régression Logistique, Random Forest, XGBoost 
    et LightGBM. Le meilleur modèle (XGBoost) atteint un ROC-AUC de \textbf{0.664} 
    sur l'ensemble de test.
  \end{tcolorbox}

  \vspace{2cm}

  \begin{tabular}{ll}
    \toprule
    \textbf{Domaine} & Transport Aérien / Prédiction \\
    \textbf{Type de tâche} & Classification Binaire Supervisée \\
    \textbf{Dataset} & 3 000 000 vols, 31 variables \\
    \textbf{Variable cible} & Retard à l'arrivée $\geq$ 15 min \\
    \textbf{Environnement} & Python 3.12, scikit-learn, XGBoost, LightGBM \\
    \textbf{Déploiement} & Application Streamlit \\
    \bottomrule
  \end{tabular}

  \vfill
  {\large\color{darkgray} Rapport d'Analyse --- 2024}
\end{titlepage}

% ─── Table des matières ───────────────────────────────────────────────────────
\tableofcontents
\newpage

% ──────────────────────────────────────────────────────────────────────────────
\section{Introduction et Contexte}
% ──────────────────────────────────────────────────────────────────────────────

\subsection{Problématique}

Les retards de vols constituent l'un des défis majeurs du transport aérien mondial.
Selon les données de l'IATA (International Air Transport Association), les retards 
coûtent annuellement plusieurs milliards d'euros aux compagnies aériennes, aux 
aéroports et aux passagers. En Europe, environ 20\% des vols arrivent avec un retard 
supérieur à 15 minutes, ce qui représente un impact économique et opérationnel considérable.

\begin{infobox}[Définition officielle du retard]
  Un vol est officiellement considéré comme \textbf{retardé} lorsque son retard 
  à l'arrivée est supérieur ou égal à \textbf{15 minutes} par rapport à l'heure 
  programmée, conformément aux normes FAA (Federal Aviation Administration) et 
  aux directives IATA.
\end{infobox}

\subsection{Objectifs}

Ce projet vise à :
\begin{itemize}[leftmargin=*]
  \item Analyser un dataset massif de vols aériens (3 millions d'entrées)
  \item Développer des modèles prédictifs pour anticiper les retards
  \item Comparer les performances de plusieurs algorithmes de classification
  \item Déployer une application interactive de prédiction en temps réel
  \item Formuler des recommandations pratiques pour les opérateurs aériens
\end{itemize}

\subsection{Enjeux Industriels}

La prédiction précoce des retards permet :
\begin{enumerate}[leftmargin=*]
  \item D'optimiser l'affectation des équipages et des aéronefs
  \item De notifier les passagers à l'avance pour réduire l'insatisfaction
  \item De minimiser les effets cascade (retards propagés entre vols)
  \item D'améliorer la gestion des portes d'embarquement dans les hubs
  \item De réduire les coûts opérationnels liés aux compensations passagers
\end{enumerate}

\newpage
% ──────────────────────────────────────────────────────────────────────────────
\section{Description des Données}
% ──────────────────────────────────────────────────────────────────────────────

\subsection{Structure du Dataset}

Le dataset \texttt{data\_vols\_1.csv} contient \textbf{3 000 000 enregistrements} 
de vols commerciaux avec 31 variables couvrant les informations de départ, d'arrivée, 
les retards par catégorie et les caractéristiques opérationnelles.

\begin{table}[H]
  \centering
  \caption{Statistiques globales du dataset}
  \begin{tabular}{lrl}
    \toprule
    \textbf{Indicateur} & \textbf{Valeur} & \textbf{Unité} \\
    \midrule
    Total d'enregistrements & 3 000 000 & vols \\
    Variables disponibles & 31 & colonnes \\
    Vols retenus (après nettoyage) & 490 011 & vols \\
    Vols retardés ($\geq$ 15 min) & 93 262 & vols \\
    Taux de retard global & 19,03 & \% \\
    Retard moyen à l'arrivée & 4,92 & minutes \\
    Retard médian à l'arrivée & $-5,0$ & minutes \\
    Retard maximum observé & 1 554 & minutes \\
    Aéroports de départ distincts & 318 & aéroports \\
    Compagnies aériennes & 12 & compagnies \\
    Types d'avions & 4 106 & aéronefs \\
    \bottomrule
  \end{tabular}
\end{table}

\subsection{Variables Principales}

\begin{table}[H]
  \centering
  \caption{Description des variables clés}
  \small
  \begin{tabular}{p{4.5cm}p{3cm}p{6cm}}
    \toprule
    \textbf{Variable} & \textbf{Type} & \textbf{Description} \\
    \midrule
    RETARD À L'ARRIVÉE & Float (cible) & Retard en minutes à l'arrivée \\
    AÉROPORT DÉPART & Catégorielle & Code IATA de l'aéroport d'origine \\
    AÉROPORT ARRIVÉE & Catégorielle & Code IATA de l'aéroport de destination \\
    COMPAGNIE AÉRIENNE & Catégorielle & Code de la compagnie opérante \\
    DÉPART PROGRAMMÉ & Entier & Heure de départ prévue (HHMM) \\
    DISTANCE & Entier & Distance du vol en kilomètres \\
    TEMPS PROGRAMMÉ & Float & Durée de vol programmée (minutes) \\
    DATE & String & Date du vol (JJ/MM/AAAA) \\
    NIVEAU DE SÉCURITÉ & Entier & Score de sécurité opérationnelle \\
    RETARD MÉTÉO & Float & Part du retard due à la météo \\
    RETARD COMPAGNIE & Float & Part du retard due à la compagnie \\
    RETARD AVION & Float & Part du retard due à l'appareil \\
    ANNULATION & Binaire & Indicateur d'annulation du vol \\
    DÉTOURNEMENT & Binaire & Indicateur de détournement \\
    \bottomrule
  \end{tabular}
\end{table}

\newpage
% ──────────────────────────────────────────────────────────────────────────────
\section{Nettoyage et Préparation des Données}
% ──────────────────────────────────────────────────────────────────────────────

\subsection{Analyse des Valeurs Manquantes}

Le dataset présente des valeurs manquantes significatives dans certaines colonnes,
notamment les causes de retard :

\begin{itemize}[leftmargin=*]
  \item \textbf{RAISON D'ANNULATION} : 98,26\% de valeurs manquantes (normale pour les vols non annulés)
  \item \textbf{RETARD SYSTÈME / SÉCURITÉ / COMPAGNIE / AVION / MÉTÉO} : 81,48\% manquants
  \item \textbf{COMPAGNIE AÉRIENNE} : 4,75\% manquants
  \item \textbf{HEURE DE DÉPART / RETARD DE DÉPART} : 1,65\% manquants
\end{itemize}

\begin{warningbox}[Décision de traitement]
  Les colonnes de causes de retard (RETARD SYSTÈME, RETARD MÉTÉO, etc.) ont été 
  exclues des features d'entraînement car elles ne sont pas disponibles \textit{a priori} 
  avant le vol --- les inclure constituerait une fuite de données (\textit{data leakage}).
\end{warningbox}

\subsection{Étapes de Nettoyage}

\begin{enumerate}[leftmargin=*]
  \item \textbf{Exclusion des vols annulés} : 1 742 vols supprimés (0,58\%)
  \item \textbf{Exclusion des vols détournés} : 297 vols supprimés (0,10\%)
  \item \textbf{Suppression des lignes sans retard} : 2 039 vols sans information d'arrivée
  \item \textbf{Création de la variable cible} : $\text{RETARDÉ} = \mathbb{1}[\text{Retard} \geq 15]$
  \item \textbf{Extraction de features temporelles} depuis la colonne DATE
\end{enumerate}

\subsection{Ingénierie des Variables (Feature Engineering)}

À partir des données brutes, les features suivantes ont été construites :

\begin{lstlisting}[caption={Création des features temporelles et encodage}]
# Variable cible binaire
df['RETARDE'] = (df["RETARD A L'ARRIVEE"] >= 15).astype(int)

# Features temporelles
df['DATE']       = pd.to_datetime(df['DATE'], dayfirst=True)
df['MOIS']       = df['DATE'].dt.month        # Saisonnalité
df['JOUR_SEMAINE'] = df['DATE'].dt.dayofweek  # Lundi=0, Dimanche=6
df['HEURE_DEP']  = df['DEPART PROGRAMME'] // 100  # Heure de départ

# Encodage des variables catégorielles
le_dep  = LabelEncoder()
le_arr  = LabelEncoder()
le_comp = LabelEncoder()
\end{lstlisting}

\subsection{Variables Utilisées pour la Modélisation}

\begin{table}[H]
  \centering
  \caption{Features retenues pour l'entraînement des modèles}
  \begin{tabular}{llp{6cm}}
    \toprule
    \textbf{Feature} & \textbf{Type} & \textbf{Justification} \\
    \midrule
    DÉPART PROGRAMMÉ & Numérique & Heure de vol intégrale \\
    HEURE\_DEP & Numérique & Heure extraite (0–23) \\
    MOIS & Numérique & Saisonnalité annuelle \\
    JOUR\_SEMAINE & Numérique & Demande hebdomadaire \\
    DISTANCE & Numérique & Longueur du trajet \\
    TEMPS PROGRAMMÉ & Numérique & Durée de vol prévue \\
    NIVEAU DE SÉCURITÉ & Numérique & Score opérationnel \\
    AÉROPORT\_DEP\_ENC & Encodé & Aéroport d'origine \\
    AÉROPORT\_ARR\_ENC & Encodé & Aéroport de destination \\
    COMPAGNIE\_ENC & Encodé & Opérateur aérien \\
    \bottomrule
  \end{tabular}
\end{table}

\textbf{Partitionnement :} 80\% entraînement (373 419 vols) / 20\% test (93 355 vols),
avec stratification sur la variable cible pour maintenir la proportion de classes.

\newpage
% ──────────────────────────────────────────────────────────────────────────────
\section{Analyse Exploratoire des Données (EDA)}
% ──────────────────────────────────────────────────────────────────────────────

\subsection{Distribution de la Variable Cible}

Sur l'ensemble des 490 011 vols retenus :
\begin{itemize}[leftmargin=*]
  \item \textbf{396 749 vols} (80,97\%) arrivent à l'heure ou avec moins de 15 minutes de retard
  \item \textbf{93 262 vols} (19,03\%) sont considérés comme retardés ($\geq$ 15 min)
\end{itemize}

Ce déséquilibre de classes (ratio 4:1) est typique des problèmes de prédiction de 
retards et a orienté nos choix de modèles et d'hyperparamètres (utilisation de 
\texttt{class\_weight='balanced'} et de \texttt{scale\_pos\_weight}).

\subsection{Analyse Temporelle}

\subsubsection{Saisonnalité Mensuelle}

L'analyse mensuelle révèle une saisonnalité marquée :
\begin{itemize}[leftmargin=*]
  \item Les mois d'été (\textbf{juin, juillet, août}) présentent des taux de retard 
        plus élevés, liés à la densité du trafic et aux conditions météorologiques
  \item \textbf{Décembre et janvier} montrent également une hausse due aux conditions 
        hivernales (verglas, neige, visibilité réduite)
  \item Les mois de \textbf{printemps} (mars--mai) sont généralement plus favorables
\end{itemize}

\subsubsection{Effet du Jour de la Semaine}

\begin{itemize}[leftmargin=*]
  \item \textbf{Vendredi et dimanche} : taux de retard les plus élevés, correspondant 
        aux pics de trafic de fin de semaine
  \item \textbf{Mardi et mercredi} : jours les plus fluides, avec moins de congestion
\end{itemize}

\subsubsection{Impact de l'Heure de Départ}

La période de départ est l'un des facteurs les plus discriminants :
\begin{itemize}[leftmargin=*]
  \item \textbf{Vols matinaux} (5h--9h) : taux de retard minimal, l'avion commence 
        sa rotation sans subir les retards accumulés de la journée
  \item \textbf{Vols du soir} (18h--22h) : taux de retard maximal, effet cascades 
        des retards précédents
  \item Tendance croissante au fil de la journée (phénomène de \textit{delay propagation})
\end{itemize}

\subsection{Analyse Géographique}

Les 318 aéroports du dataset présentent des taux de retard très hétérogènes.
Les aéroports avec les taux les plus élevés sont généralement :
\begin{itemize}[leftmargin=*]
  \item Des hubs importants avec une forte densité de trafic
  \item Des aéroports exposés à des conditions météorologiques défavorables
  \item Des plateformes de correspondance subissant l'effet cascade
\end{itemize}

\subsection{Distribution du Retard en Minutes}

La distribution du retard à l'arrivée présente une forme \textbf{asymétrique à droite} 
(skewness positif) :
\begin{itemize}[leftmargin=*]
  \item Médiane : $-5$ minutes (la majorité des vols arrivent légèrement en avance)
  \item Moyenne : $+4{,}92$ minutes (tirée vers le haut par les retards extrêmes)
  \item Maximum observé : $1\,554$ minutes ($\approx 26$ heures)
  \item Les valeurs extrêmes correspondent à des événements exceptionnels (pannes majeures, 
        conditions climatiques sévères)
\end{itemize}

\newpage
% ──────────────────────────────────────────────────────────────────────────────
\section{Modèles de Machine Learning}
% ──────────────────────────────────────────────────────────────────────────────

\subsection{Choix des Algorithmes}

Quatre algorithmes ont été sélectionnés pour couvrir un spectre de complexité :

\begin{infobox}[Stratégie de sélection des modèles]
  Le choix s'est porté sur : (1) un modèle baseline linéaire, (2) un modèle 
  d'ensemble basé sur les arbres, et (3) deux algorithmes de gradient boosting 
  qui représentent l'état de l'art pour les données tabulaires.
\end{infobox}

\subsubsection{1. Régression Logistique (Baseline)}

Modèle linéaire servant de référence (\textit{baseline}). Avantages : 
interprétabilité maximale, rapidité d'entraînement. Limites : ne capture pas 
les interactions non-linéaires entre variables.

\textbf{Hyperparamètres :} \texttt{max\_iter=500}, \texttt{class\_weight='balanced'}

\subsubsection{2. Random Forest}

Ensemble de 100 arbres de décision entraînés en parallèle. 
Robuste au bruit, peu sensible au réglage des hyperparamètres.
Bonne interprétabilité via l'importance des features.

\textbf{Hyperparamètres :} \texttt{n\_estimators=100}, \texttt{max\_depth=12}, 
\texttt{class\_weight='balanced'}

\subsubsection{3. XGBoost (Extreme Gradient Boosting)}

Algorithme de boosting par gradient utilisant des arbres CART. 
Régularisation L1/L2 intégrée, gestion native des valeurs manquantes.
Optimal pour les données tabulaires structurées.

\textbf{Hyperparamètres :} \texttt{n\_estimators=200}, \texttt{max\_depth=6}, 
\texttt{learning\_rate=0.1}, \texttt{scale\_pos\_weight=4.25}

\subsubsection{4. LightGBM (Light Gradient Boosting Machine)}

Variante de gradient boosting utilisant une stratégie de croissance feuille 
par feuille (\textit{leaf-wise}) plutôt que niveau par niveau. 
Significativement plus rapide que XGBoost sur de grands datasets.

\textbf{Hyperparamètres :} \texttt{n\_estimators=200}, \texttt{max\_depth=8}, 
\texttt{learning\_rate=0.1}, \texttt{class\_weight='balanced'}

\subsection{Gestion du Déséquilibre de Classes}

Le ratio 80\%/20\% entre vols non-retardés et retardés nécessite des stratégies adaptées :

\begin{table}[H]
  \centering
  \caption{Stratégies de gestion du déséquilibre par modèle}
  \begin{tabular}{ll}
    \toprule
    \textbf{Modèle} & \textbf{Stratégie} \\
    \midrule
    Régression Logistique & \texttt{class\_weight='balanced'} \\
    Random Forest & \texttt{class\_weight='balanced'} \\
    XGBoost & \texttt{scale\_pos\_weight = n\_neg / n\_pos $\approx$ 4.25} \\
    LightGBM & \texttt{class\_weight='balanced'} \\
    \bottomrule
  \end{tabular}
\end{table}

\newpage
% ──────────────────────────────────────────────────────────────────────────────
\section{Évaluation des Performances}
% ──────────────────────────────────────────────────────────────────────────────

\subsection{Métriques d'Évaluation}

Pour un problème de classification avec déséquilibre de classes, plusieurs métriques 
complémentaires sont nécessaires :

\begin{align}
  \text{Accuracy} &= \frac{VP + VN}{VP + VN + FP + FN} \\[6pt]
  \text{Précision} &= \frac{VP}{VP + FP} \\[6pt]
  \text{Rappel} &= \frac{VP}{VP + FN} \\[6pt]
  \text{F1-Score} &= 2 \cdot \frac{\text{Précision} \times \text{Rappel}}{\text{Précision} + \text{Rappel}} \\[6pt]
  \text{ROC-AUC} &= \int_0^1 \text{TPR}(\text{FPR})\, d(\text{FPR})
\end{align}

où VP = Vrais Positifs, VN = Vrais Négatifs, FP = Faux Positifs, FN = Faux Négatifs.

\subsection{Tableau Comparatif des Résultats}

\begin{table}[H]
  \centering
  \caption{Comparaison des performances sur l'ensemble de test (93 355 vols)}
  \setlength{\tabcolsep}{8pt}
  \begin{tabular}{lcccccc}
    \toprule
    \textbf{Modèle} & \textbf{Accuracy} & \textbf{F1-Score} & \textbf{Précision} & \textbf{Rappel} & \textbf{ROC-AUC} & \textbf{Avg. Prec.} \\
    \midrule
    Rég. Logistique  & 0.5671 & 0.3381 & 0.2385 & 0.5801 & 0.5964 & 0.2387 \\
    Random Forest    & 0.6241 & 0.3673 & 0.2704 & 0.5726 & 0.6476 & 0.2986 \\
    \rowcolor{best}
    \textbf{XGBoost} & \textbf{0.6176} & \textbf{0.3820} & \textbf{0.2760} & \textbf{0.6202} & \textbf{0.6644} & \textbf{0.3153} \\
    LightGBM         & 0.6128 & 0.3801 & 0.2735 & 0.6230 & 0.6638 & 0.3150 \\
    \bottomrule
  \end{tabular}
  \\\vspace{4pt}
  \small \textit{La ligne en vert correspond au meilleur modèle selon le ROC-AUC}
\end{table}

\subsection{Analyse des Résultats par Modèle}

\subsubsection{Régression Logistique}

Le modèle baseline produit les performances les plus faibles mais reste au-dessus 
d'un classificateur aléatoire (ROC-AUC = 0.50). Son accuracy de 56,7\% avec 
un rappel de 58\% montre qu'il détecte une fraction raisonnable des retards,
mais génère de nombreux faux positifs (précision 23,9\%).

\subsubsection{Random Forest}

Nette amélioration par rapport au baseline : ROC-AUC passe à 0.648 (+8,6\%).
L'accuracy de 62,4\% est la plus élevée de tous les modèles, mais cela s'explique
en partie par un rappel légèrement plus bas (57,3\%). Le Random Forest offre une 
bonne stabilité grâce à l'agrégation de 100 arbres.

\subsubsection{XGBoost --- Meilleur Modèle}

\begin{resultbox}[XGBoost : Modèle Recommandé]
  XGBoost obtient le \textbf{meilleur ROC-AUC} (0.664) et la meilleure 
  \textbf{Average Precision} (0.315), ce qui en fait le modèle le plus 
  performant pour discriminer les vols retardés des vols à l'heure. 
  Son rappel de 62\% signifie qu'il détecte 6 retards sur 10 correctement.
\end{resultbox}

\subsubsection{LightGBM}

Performances quasi-identiques à XGBoost (ROC-AUC 0.664 vs 0.663), avec l'avantage 
d'un temps d'entraînement significativement réduit sur de grands datasets. 
En production, LightGBM serait préféré pour sa scalabilité.

\subsection{Interprétation des Matrices de Confusion}

Pour XGBoost sur 93 355 vols de test :
\begin{itemize}[leftmargin=*]
  \item \textbf{Vrais Positifs} : $\approx$ 11 500 retards correctement prédits
  \item \textbf{Faux Négatifs} : $\approx$ 7 000 retards manqués (vols annoncés à l'heure, finalement retardés)
  \item \textbf{Faux Positifs} : $\approx$ 30 000 fausses alarmes
  \item \textbf{Vrais Négatifs} : $\approx$ 44 000 vols à l'heure correctement identifiés
\end{itemize}

Le ratio FP élevé est inévitable avec un déséquilibre 80/20 et des features 
incomplètes (absence des données météo temps-réel, état des pistes, etc.).

\subsection{Importance des Variables}

L'analyse d'importance des features (Random Forest, XGBoost, LightGBM) révèle
un consensus sur les variables les plus prédictives :

\begin{enumerate}[leftmargin=*]
  \item \textbf{Distance du vol} : les longs-courriers accumulent plus facilement des retards
  \item \textbf{Temps de vol programmé} : corrélé à la distance, proxy de la complexité
  \item \textbf{Heure de départ} : effet cascade quotidien
  \item \textbf{Aéroport de départ} : certains hubs sont structurellement plus congestionné
  \item \textbf{Compagnie aérienne} : reflète les pratiques opérationnelles différentes
  \item \textbf{Mois} : saisonnalité et trafic
  \item \textbf{Jour de la semaine} : pattern hebdomadaire de demande
\end{enumerate}

\newpage
% ──────────────────────────────────────────────────────────────────────────────
\section{Limites et Améliorations Proposées}
% ──────────────────────────────────────────────────────────────────────────────

\subsection{Limites Actuelles}

\begin{warningbox}[Limites identifiées]
  \begin{itemize}[leftmargin=*, nosep]
    \item \textbf{Absence de données météo temps-réel} : vent, visibilité, précipitations sont 
          les principaux facteurs non capturés
    \item \textbf{Absence d'historique de l'avion} : le retard précédent de l'appareil 
          (rotation) est très prédictif mais non disponible
    \item \textbf{Déséquilibre de classes} : malgré les stratégies de pondération, 
          le modèle reste biaisé
    \item \textbf{Features statiques} : les features capturent des tendances globales, 
          pas les conditions dynamiques du jour J
    \item \textbf{Échantillonnage} : 500 000 vols utilisés sur 3 millions disponibles
  \end{itemize}
\end{warningbox}

\subsection{Pistes d'Amélioration}

\subsubsection{Amélioration des Données}

\begin{itemize}[leftmargin=*]
  \item \textbf{Intégration météo} : API OpenWeatherMap ou Météo-France pour données 
        historiques par aéroport et date
  \item \textbf{Données de rotation} : retard du vol précédent effectué par le même avion
  \item \textbf{Données NAS (National Airspace System)} : saturation du contrôle aérien
  \item \textbf{Données aéroport} : taux d'occupation des pistes, incidents récents
  \item \textbf{Features cycliques} : encoder l'heure et le mois comme $\sin(2\pi x / T)$, 
        $\cos(2\pi x / T)$ pour préserver la continuité
\end{itemize}

\subsubsection{Amélioration des Modèles}

\begin{itemize}[leftmargin=*]
  \item \textbf{Optimisation des hyperparamètres} : recherche bayésienne (Optuna) 
        plutôt que manuelle
  \item \textbf{SMOTE / ADASYN} : suréchantillonnage synthétique de la classe minoritaire
  \item \textbf{Modèles neuronaux} : TabNet, réseaux de neurones tabulaires
  \item \textbf{Threshold tuning} : optimiser le seuil de décision (par défaut 0.5) 
        pour maximiser le rappel ou le F1
  \item \textbf{Stacking / Blending} : combiner les prédictions des 4 modèles
  \item \textbf{Entraînement sur dataset complet} : utiliser les 3 millions de vols
\end{itemize}

\subsubsection{Amélioration du Déploiement}

\begin{itemize}[leftmargin=*]
  \item \textbf{Mise à jour incrémentale} : réentraînement périodique avec les nouveaux vols
  \item \textbf{API REST} : exposer le modèle via FastAPI pour intégration systèmes
  \item \textbf{Monitoring} : suivi de la dérive des données (data drift) en production
  \item \textbf{A/B Testing} : comparer les décisions prises avec/sans le système
  \item \textbf{Explicabilité SHAP} : expliquer chaque prédiction individuellement
\end{itemize}

\newpage
% ──────────────────────────────────────────────────────────────────────────────
\section{Déploiement de l'Application Streamlit}
% ──────────────────────────────────────────────────────────────────────────────

\subsection{Architecture de l'Application}

L'application Streamlit offre une interface utilisateur intuitive permettant :
\begin{itemize}[leftmargin=*]
  \item La prédiction interactive pour un vol individuel
  \item La visualisation des performances des 4 modèles
  \item L'exploration des analyses EDA
  \item La comparaison des prédictions entre modèles
\end{itemize}

\subsection{Structure du Projet}

\begin{lstlisting}[language=bash, caption={Structure des fichiers du projet}]
flight_delay_project/
|-- src/
|   |-- app.py              # Application principale Streamlit
|-- outputs/
|   |-- eda_analysis.png    # Figures EDA
|   |-- model_comparison.png
|   |-- confusion_matrices.png
|   |-- feature_importance.png
|   |-- roc_curves.png
|   |-- results_table.json
|   |-- stats.json
|-- models/
|   |-- regression_logistique.pkl
|   |-- random_forest.pkl
|   |-- xgboost.pkl
|   |-- lightgbm.pkl
|   |-- le_dep.pkl
|   |-- le_arr.pkl
|   |-- le_comp.pkl
|-- reports/
    |-- rapport_retards_vols.tex
\end{lstlisting}

\subsection{Instructions de Déploiement}

\begin{lstlisting}[language=bash, caption={Commandes de déploiement}]
# Installation des dependances
pip install streamlit scikit-learn xgboost lightgbm
pip install pandas numpy matplotlib seaborn joblib

# Lancement de l'application
cd flight_delay_project/
streamlit run src/app.py --server.port 8501

# Acces dans le navigateur
# http://localhost:8501
\end{lstlisting}

\subsection{Fonctionnalités de l'Application}

\begin{table}[H]
  \centering
  \caption{Pages et fonctionnalités de l'application Streamlit}
  \begin{tabular}{p{3.5cm}p{9cm}}
    \toprule
    \textbf{Page} & \textbf{Fonctionnalités} \\
    \midrule
    Accueil & Métriques clés, description du projet, résumé du meilleur modèle \\
    Prédiction & Formulaire interactif, résultat avec probabilités, comparaison entre modèles \\
    Performance Modèles & Tableau comparatif avec surbrillance du meilleur, courbes ROC, matrices de confusion, importance des features \\
    Analyse EDA & Visualisations de la distribution des retards, tendances temporelles, géographie \\
    \bottomrule
  \end{tabular}
\end{table}

\newpage
% ──────────────────────────────────────────────────────────────────────────────
\section{Conclusion}
% ──────────────────────────────────────────────────────────────────────────────

\subsection{Synthèse}

Ce projet a démontré la faisabilité d'un système de prédiction de retards de vols 
basé sur le Machine Learning, appliqué à un dataset industriel de grande taille 
(3 millions de vols). 

Les principaux résultats sont :

\begin{resultbox}[Résultats Clés]
  \begin{itemize}[leftmargin=*, nosep]
    \item \textbf{XGBoost} est le meilleur modèle avec un ROC-AUC de \textbf{0.664}
    \item \textbf{LightGBM} offre des performances quasi-équivalentes avec un avantage en vitesse
    \item Les features temporelles (heure, mois) et la distance sont les plus prédictives
    \item Un taux de rappel de 62\% permet de détecter la majorité des retards significatifs
    \item L'application Streamlit déployée rend le modèle accessible à des non-experts
  \end{itemize}
\end{resultbox}

\subsection{Perspectives}

L'intégration de données météorologiques temps-réel et de l'historique de rotation 
des avions constitue la piste d'amélioration la plus prometteuse. Des études ont 
montré que ces deux types de données peuvent porter le ROC-AUC au-delà de 0.80 
sur des problèmes similaires.

Pour un déploiement industriel, les prochaines étapes seraient :
\begin{enumerate}[leftmargin=*]
  \item Développer une API REST pour l'intégration dans les systèmes de gestion aéroportuaire
  \item Mettre en place un pipeline de réentraînement automatique (MLOps)
  \item Implémenter un système d'explication des prédictions (SHAP values)
  \item Réaliser un test en conditions réelles sur un aéroport pilote
  \item Évaluer l'impact économique des prédictions correctes vs erreurs de prédiction
\end{enumerate}

\subsection{Valeur Métier}

Un système de prédiction performant permet d'économiser en moyenne entre 
\textbf{500 et 2 000 euros par vol retardé} identifié à l'avance, grâce à :
\begin{itemize}[leftmargin=*]
  \item La notification proactive des passagers (réduction des compensations)
  \item L'optimisation du routing des équipages et des avions
  \item La réduction des effets cascade sur les vols suivants
  \item L'amélioration de la satisfaction client et de l'image de marque
\end{itemize}

Appliqué aux 93 262 vols retardés de notre dataset (19\% du total), 
un système avec 62\% de rappel permettrait de traiter proactivement 
environ \textbf{57 822 retards}, représentant un potentiel d'économies 
de plusieurs dizaines de millions d'euros sur l'ensemble du réseau.

\vspace{1cm}
\begin{center}
  \rule{0.5\textwidth}{0.4pt}\\[0.3cm]
  {\small\itshape Rapport généré dans le cadre d'un projet d'analyse de données --- 2024}\\
  {\small\itshape Technologies : Python 3.12 $\cdot$ scikit-learn $\cdot$ XGBoost $\cdot$ LightGBM $\cdot$ Streamlit}
\end{center}

\end{document}
